\subsection{Gesamtbewertung}
\label{subsec:gesamtbewertung}

Tabelle~\ref{tab:template-evaluation} unterscheidet weitgehend und teilweise
adressierte Kriterien von nicht abschließend nachgewiesenen. Diese Stufen sind
qualitativ und messen weder Qualität noch Produktivität numerisch.

\begin{table}[htbp]
  \centering
  \scriptsize
  \renewcommand{\arraystretch}{1.08}
  \caption{Zusammengeführte Bewertung des Technologie-Templates}
  \label{tab:template-evaluation}
  \begin{tabularx}{\textwidth}{@{}>{\raggedright\arraybackslash}p{0.18\textwidth}>{\raggedright\arraybackslash}p{0.31\textwidth}>{\raggedright\arraybackslash}p{0.19\textwidth}>{\raggedright\arraybackslash}X@{}}
    \toprule
    \textbf{Kriterium} & \textbf{Evidenzbasis} & \textbf{Bewertung} &
      \textbf{wesentliche Einschränkung} \\
    \midrule
    Wiederverwendung und Standardisierung & Gemeinsamer Kern, fünf Profile,
      Generator und einheitliche Befehlsoberfläche & weitgehend adressiert &
      Gilt für den definierten Stack; keine automatische Aktualisierung
      abgeleiteter Projekte. \\
    Variabilität und Erweiterbarkeit & Profil-Presets, Capability-Abhängigkeiten
      und geprüfte PostgreSQL-Kombinationen & weitgehend adressiert & Kleiner
      aktueller Katalog; Erweiterungen vergrößern Pflege- und Testumfang. \\
    Testbarkeit und Reproduzierbarkeit & Dokumentierte Profil-, Integrations-
      und Buildabnahme sowie ein byte-identischer Generierungsvergleich &
      teilweise adressiert & Aktuelle Remote-Jobs teilweise fehlgeschlagen;
      Reproduzierbarkeitsversuch nur für eine Kombination. \\
    Entwicklerfreundlichkeit und Dokumentierbarkeit & \texttt{control.py},
      Diagnosewege, versionierte Manifeste und gegliederte Dokumentation &
      teilweise adressiert & Onboarding und Bedienaufwand nicht empirisch
      gemessen; mehrere Toolchains bleiben sichtbar. \\
    Wartbarkeit & Eine zentrale Baseline mit gemeinsamen Tests und
      Dokumentationsregeln & teilweise adressiert & Technologievielfalt,
      Variantenregression und manuelle Schnittstellensynchronisierung. \\
    Codequalitäts- und Architektur-Governance & Zentrale Policy, 13 Rule IDs,
      135 fokussierte Tests, lokale Befehle und Master-Quality-Job &
      teilweise adressiert & Generierte Backendprofile nicht grün;
      vollständiger lokaler Lauf durch Rust-Systemabhängigkeiten begrenzt;
      \texttt{CQ001} unterdrückbar und kein Clean-Code-Nachweis. \\
    Plattformabdeckung und Betriebsreife & Web-, Backend-, PostgreSQL-,
      Container-Image-, Linux- und macOS-Nachweise; providerneutrale Grenze &
      nicht abschließend nachgewiesen & Windows-Paket, Compose-Start, Signing,
      Notarisierung und produktives Deployment nicht belegt. \\
    \bottomrule
  \end{tabularx}
  \smallskip
  {\footnotesize Quelle: Eigene Bewertung auf Grundlage der vorangehenden
  Kapitel, \textcite{kleiveist2026acceptance} und der in
  Abschnitt~\ref{subsec:continuous-integration} ausgewerteten CI-Läufe.\par}
\end{table}


Die Baseline ist für die definierten Profile technisch tragfähig. Kern,
Presets, Abhängigkeiten und Tooling unterstützen Wiederverwendung,
Standardisierung und Flexibilität. Die Abnahme stützt Generierung,
Installation, Tests, Builds, PostgreSQL-Kombinationen und Linux-Packaging
\parencite{kleiveist2026acceptance}. Neuere externe Läufe ergänzen Web-,
Container-, Linux- und macOS-Nachweise, enthalten aber Fehler in Tooling-,
Desktopprofil- und Windows-Jobs.
Die Evidenz ist damit belastbar, aber nach Commit und Plattform unterschiedlich
weitreichend.

Grenzen sind zentrale Pflegekomplexität, manuell synchronisierte Verträge und
fehlende Updates generierter Projekte. Windows-Paketierung, signierte Releases,
Notarisierung und produktiver Cloudbetrieb sind nicht nachgewiesen. Bewusst
ausgelassene Provider- und Fachlogik ist dagegen kein Fehler.
Tabelle~\ref{tab:template-evaluation} macht den Unterschied zwischen
adressierter Anforderung und offener Nachweisbreite sichtbar.

Automatisierung und gemeinsame Workflows schaffen plausibles, aber nicht
quantifiziertes Produktivitätspotenzial. Der Stand eignet sich als
standardisierte Basis für die vorgesehenen Projektfamilien, nicht als fertige
Produktionsanwendung oder universelle Vorlage.
Ob sich die zentrale Investition wirtschaftlich lohnt, bleibt ohne Nutzungs-
und Kostendaten offen.
\beginnerinput{workspace/statement/700_bewertung/760}
